% =====================================================================
\chapter{Mezcla de expertos}
\label{tema:4}
% =====================================================================

Un Transformer denso aplica a cada token todos sus parámetros. Una capa de \emph{mezcla de expertos} (MoE) sustituye el MLP por varios MLP alternativos y por un enrutador que elige unos pocos para cada token, de modo que los parámetros crecen con el número de expertos y el cómputo por token solo con el número de expertos elegidos. Este tema formula la capa, distingue parámetros totales y activos, estudia el problema del balanceo de carga y sus soluciones, y reconstruye los recuentos publicados de varios modelos. Se sigue la clase~4 de CS336~\cite{cs336}, los artículos de Shazeer \emph{et al.}~\cite[§2 y §4]{shazeer2017}, Fedus \emph{et al.}~\cite[§2]{switch} y Jiang \emph{et al.}~\cite[§2]{mixtral}, y los informes de DeepSeek-V2 y DeepSeek-V3~\cite[§2.2]{deepseekv2},~\cite[§2.1.2]{deepseekv3}.

\section{La capa de mezcla de expertos}

\begin{definicion}[Capa MoE con enrutado top-$k$]\label{def:moe}
Sean $E_1,\dots,E_N:\R^{1\times d}\to\R^{1\times d}$ MLP con la misma forma (los \emph{expertos}), $W_r\in\R^{d\times N}$ la matriz del \emph{enrutador} y $1\le k\le N$. Para un token $x$, sean $z=xW_r\in\R^{N}$ sus logits y $\mathcal T(x)$ el conjunto de los índices de los $k$ logits mayores. La salida de la capa es
\[
y=\sum_{i\in\mathcal T(x)}g_i(x)\,E_i(x),\qquad g_i(x)=\frac{e^{z_i}}{\sum_{j\in\mathcal T(x)}e^{z_j}} .
\]
Los expertos que no están en $\mathcal T(x)$ no se evalúan. Esta es la formulación de Mixtral, $G(x)=\softmax(\TopK(xW_r))$~\cite[§2.1]{mixtral}.
\end{definicion}

\addtocounter{figure}{1}%

Hay variantes en el orden de las operaciones. Shazeer \emph{et al.} también normalizan solo entre los elegidos, pero añaden ruido gaussiano a los logits, con una desviación aprendida, $H(x)_i=(xW_r)_i+\xi_i\,\softplus\bigl((xW_{\mathrm{ruido}})_i\bigr)$ con $\xi_i\sim\mathcal N(0,1)$, para favorecer la exploración y el balanceo~\cite[§2.1]{shazeer2017}. El Switch Transformer usa $k=1$ y pondera la salida por la probabilidad de la softmax sobre los $N$ expertos, sin renormalizar~\cite[§2.1]{switch}. La clase~4 de CS336 recoge que DeepSeek-V1 y V2, Grok y Qwen aplican la softmax antes del top-$k$, y que Mixtral, DBRX y DeepSeek-V3 normalizan después~\cite{cs336}.

\begin{proposicion}[Renormalización]\label{prop:renorm}
Sean $p=\softmax(z)$ las probabilidades sobre los $N$ expertos y $g$ los pesos de la \cref{def:moe}. Para $i\in\mathcal T(x)$, $g_i=p_i/\sum_{j\in\mathcal T(x)}p_j\ge p_i$, con igualdad para todo $i$ si y solo si $k=N$. Además, el orden de los expertos elegidos es el mismo con $g$ que con $p$.
\end{proposicion}

\addtocounter{definicion}{1}%

\section{Parámetros totales y activos}

\begin{proposicion}[Parámetros totales y activos]\label{prop:activos}
Sea un modelo con $L$ capas MoE de $N$ expertos de $P_E$ parámetros cada uno, $k$ de ellos activos por token, y sea $P_0$ el número de parámetros compartidos por todos los tokens (atención, \emph{embeddings}, normalizaciones y enrutadores). Entonces
\[
\Ntot=P_0+LNP_E,\qquad \Nact=P_0+LkP_E .
\]
El cómputo por token en la pasada hacia delante es de unos $2\Nact$ FLOPs y el de entrenamiento con $D$ tokens, $C\approx 6\Nact D$, mientras que la memoria de los pesos es proporcional a $\Ntot$.
\end{proposicion}

\addtocounter{definicion}{1}%

\section{Balanceo de carga}

Si un experto empieza siendo algo mejor, el enrutador le envía más tokens, recibe más gradiente, mejora más y acaba acaparando el tráfico; los demás apenas se entrenan. Shazeer \emph{et al.} ya observaron este colapso y lo combatieron con una pérdida sobre el coeficiente de variación de la importancia de cada experto, $\mathrm{Importance}(X)=\sum_{x\in X}G(x)$~\cite[§4]{shazeer2017}. La formulación más usada es la del Switch Transformer.

\begin{definicion}[Pérdida de balanceo]\label{def:bal}
Sea un lote $\mathcal B$ de tokens. Para cada experto $i$, sean $f_i$ la fracción de las asignaciones del lote que recibe y $P_i=\frac{1}{\abs{\mathcal B}}\sum_{x\in\mathcal B}p_i(x)$ la probabilidad media que le da el enrutador. La \emph{pérdida de balanceo} es
\[
\Lbal=N\sum_{i=1}^{N}f_iP_i,
\]
y se suma a la entropía cruzada con un coeficiente $\alpha$; el Switch Transformer usa $\alpha=10^{-2}$~\cite[§2.2, ecs.~4--6]{switch}.
\end{definicion}

Las fracciones $f_i$ proceden de una selección discreta y no son derivables; el gradiente fluye solo a través de $P_i$. La clase~4 de CS336 lo resume así: la derivada respecto de $p_i(x)$ es proporcional al número de tokens asignados al experto~$i$, de modo que un experto más usado ve reducida su probabilidad con más fuerza~\cite{cs336}.

\begin{teorema}[Cotas de la pérdida de balanceo]\label{teo:bal}
Si $f=P$ (el tráfico real coincide con la probabilidad media), entonces $1\le\Lbal\le N$. La cota inferior se alcanza si y solo si $P_i=1/N$ para todo $i$, y la superior, si y solo si un experto recibe todo el tráfico.
\end{teorema}

\begin{table}[ht]
\centering\small
\caption{Valores de $\Lbal$ con $N=4$.}\label{tab:bal}
\begin{tabular}{@{}lllc@{}}
\toprule
Situación & $f$ & $P$ & $\Lbal$ \\
\midrule
Reparto uniforme & $(\frac14,\frac14,\frac14,\frac14)$ & $(\frac14,\frac14,\frac14,\frac14)$ & $1$ \\
Desequilibrio & $(0{,}5;\ 0{,}3;\ 0{,}1;\ 0{,}1)$ & $(0{,}4;\ 0{,}3;\ 0{,}2;\ 0{,}1)$ & $1{,}28$ \\
Colapso & $(1,0,0,0)$ & $(1,0,0,0)$ & $4$ \\
\bottomrule
\end{tabular}
\end{table}

\begin{observacion}[Balanceo sin pérdida auxiliar]
DeepSeek-V3 sustituye casi por completo la pérdida auxiliar por un sesgo $b_i$ por experto que solo interviene en la selección: los expertos se eligen según $s_i(x)+b_i$, pero se ponderan con $s_i(x)$, y al final de cada paso $b_i$ disminuye en una cantidad fija $\gamma$ si el experto ha estado sobrecargado y aumenta en $\gamma$ si ha estado infrautilizado~\cite[§2.1.2]{deepseekv3}. Mantiene, con un coeficiente muy pequeño, una pérdida de balanceo por secuencia. DeepSeek-V3 usa además afinidades sigmoides, $s_i(x)=\sigma(u\,e_i\T)$, en lugar de una softmax.
\end{observacion}

\section{Capacidad y descarte de tokens}

En un acelerador, cada experto procesa un tensor de tamaño fijo, de modo que el número de tokens que puede recibir por lote está acotado.

\begin{definicion}[Capacidad de un experto]\label{def:capacidad}
Con un lote de $B_{\mathrm{t}}$ tokens, $N$ expertos, $k$ asignaciones por token y un \emph{factor de capacidad} $c\ge 1$, la \emph{capacidad} de cada experto es
\[
C_{\mathrm{exp}}=c\cdot\frac{k\,B_{\mathrm{t}}}{N}
\]
tokens~\cite[§2.2, ec.~3]{switch}. Las asignaciones que exceden la capacidad se \emph{descartan}: para ellas, ese experto no se evalúa.
\end{definicion}

Un token descartado en todos sus expertos sale de la capa con su vector intacto, porque el bloque calcula $h+\FFN(\cdot)$ y la contribución del MLP es nula: no se pierde la información del token, sino la transformación de esa capa. Con $B_{\mathrm{t}}=1024$, $N=8$, $k=1$ y $c=1{,}25$, la capacidad es $1{,}25\cdot 128=160$ tokens por experto. Un factor mayor descarta menos tokens, pero reserva más memoria y cómputo para posiciones de relleno. La clase~4 de CS336 señala una consecuencia práctica: como el descarte se decide por lote, las peticiones de otros usuarios pueden provocar el descarte de los tokens propios, lo que añade no determinismo a la inferencia~\cite{cs336}.

\section{Expertos de grano fino y compartidos}

DeepSeekMoE divide cada experto en $m$ expertos más pequeños y activa $m$ veces más, de modo que los parámetros y el cómputo no cambian, y reserva algunos expertos \emph{compartidos} que procesan todos los tokens~\cite[§2.2.1]{deepseekv2}. DeepSeek-V2 usa 2 expertos compartidos y 160 enrutados, de los que activa~6; DeepSeek-V3, 1 compartido y 256 enrutados, de los que activa~8~\cite{deepseekv2,deepseekv3}.

\begin{proposicion}[Combinaciones de expertos]\label{prop:combinaciones}
Si cada uno de $N$ expertos se divide en $m$ expertos de tamaño $1/m$ y se activan $mk$ de ellos, el número de combinaciones posibles pasa de $\binom{N}{k}$ a $\binom{mN}{mk}$, con los mismos parámetros totales y activos.
\end{proposicion}

Con $N=8$, $k=2$ y $m=4$, se pasa de $\binom82=28$ combinaciones a $\binom{32}{8}=10\,518\,300$. La clase~4 de CS336 recoge que las ablaciones de DeepSeek muestran ganancias de los expertos de grano fino y de los compartidos, mientras que las de OLMoE confirman las primeras pero no las segundas~\cite{cs336}.

\section{Casos públicos}

\begin{table}[ht]
\centering\small
\caption{Modelos MoE con información publicada por sus autores.}\label{tab:moe}
\begin{tabularx}{\textwidth}{@{}l>{\raggedright\arraybackslash}X>{\raggedright\arraybackslash}X@{}}
\toprule
Modelo & Expertos y enrutado & Parámetros \\
\midrule
Shazeer \emph{et al.} (2017)~\cite{shazeer2017} & hasta $131\,072$ expertos entre capas LSTM, top-$k$ con ruido & hasta $1{,}37\cdot 10^{11}$ en la capa MoE \\
Switch-C (2021)~\cite{switch} & $2048$ expertos, top-1 & $1{,}571\cdot 10^{12}$ (tabla~9) \\
Mixtral 8x7B (2024)~\cite{mixtral} & 8 expertos, top-2 & $4{,}7\cdot 10^{10}$ totales, $1{,}3\cdot 10^{10}$ activos \\
Grok-1 (2024)~\cite{grok1} & 8 expertos, 2 por token & $3{,}14\cdot 10^{11}$ \\
DeepSeek-V2 (2024)~\cite{deepseekv2} & 2 compartidos y 160 enrutados, 6 activos & $2{,}36\cdot 10^{11}$ totales, $2{,}1\cdot 10^{10}$ activos \\
DeepSeek-V3 (2024)~\cite{deepseekv3} & 1 compartido y 256 enrutados, 8 activos & $6{,}71\cdot 10^{11}$ totales, $3{,}7\cdot 10^{10}$ activos \\
Gemini 1.5 Pro (2024)~\cite{gemini15} & «Transformer de mezcla de expertos dispersa» & no publicado \\
\bottomrule
\end{tabularx}
\end{table}

\begin{observacion}[Lo que no es público]
El informe de Gemini~1.5 describe el modelo como un Transformer MoE disperso, pero no da el número de expertos ni el de parámetros~\cite[§3.1]{gemini15}. Anthropic no ha publicado la arquitectura de Claude, y xAI no ha publicado la de las versiones de Grok posteriores a Grok-1. Nada de lo que se dice en este tema debe atribuirse a esos modelos.
\end{observacion}

\begin{observacion}[Paralelismo de expertos]
Las MoE añaden una forma de paralelismo a las del módulo~17: cada dispositivo guarda un subconjunto de expertos y los tokens viajan al dispositivo de su experto y vuelven mediante una comunicación de todos con todos (\emph{all-to-all}). DeepSeek-V2 limita el número de dispositivos a los que puede ir cada token para acotar esa comunicación~\cite[§2.2.2]{deepseekv2}.
\end{observacion}

